\label{CH:Giris}

İnternet üzerinden gerçekleştirilen e-posta iletişimi, kişisel ve kurumsal yaşamın ayrılmaz bir parçası haline gelmiştir. Ancak standart e-posta protokolleri (SMTP, IMAP, POP3) güvenlik mekanizmaları açısından oldukça yetersiz kalmaktadır: iletiler açık metin olarak aktarılabilmekte, gönderici kimliği kolayca taklit edilebilmekte ve mesaj içerikleri değiştirilebilmektedir \cite{StallingsCNS}. Bu durum; gizlilik, bütünlük, kimlik doğrulama ve inkar edememe gibi temel güvenlik gereksinimlerinin karşılanmasını zorunlu kılmaktadır.

Kriptografi, bu gereksinimleri karşılayan matematiksel temel mekanizmaları sunmaktadır. Simetrik şifreleme algoritmaları (AES), yüksek performanslı gizlilik sağlarken; asimetrik algoritmalar (RSA) güvenli anahtar değişimi ve dijital imzaya olanak tanımaktadır. Özet (hash) fonksiyonları ise mesaj bütünlüğünün denetlenmesinde kullanılmaktadır. Bu üç mekanizmanın birleşimi olan hibrit şifreleme, günümüz güvenli iletişim sistemlerinin temelini oluşturmaktadır \cite{KatzLindell}.

Bu bitirme projesi kapsamında; SHA-256 özet fonksiyonu, RSA-2048 PSS dijital imzası ve AES-256-GCM kimlik doğrulamalı şifreleme algoritmaları bir araya getirilerek e-posta mesajlarında gizlilik, bütünlük, kimlik doğrulama ve inkar edememe özelliklerini sağlayan eğitici bir hibrit şifreleme uygulaması geliştirilmiştir. Geliştirilen PyQt6 tabanlı masaüstü uygulama, Alice (gönderici) ile Bob (alıcı) senaryosu üzerinden kriptografik iş akışını adım adım görselleştirmekte ve her algoritmanın iç işleyişini animasyonlu pencereler aracılığıyla pedagojik biçimde sunmaktadır.

\textbf{Literatür Özeti}

E-posta güvenliği alanında yapılan çalışmalar, farklı kriptografik yaklaşımları ele almaktadır. PGP (Pretty Good Privacy), 1991 yılında Phil Zimmermann tarafından geliştirilen ve RSA ile IDEA algoritmalarını birleştiren ilk yaygın e-posta şifreleme standardıdır \cite{ZimmermanPGP}. OpenPGP (RFC 4880) bu standardın güncellenmiş ve açık kaynaklı hâlidir. S/MIME (RFC 5751) ise X.509 sertifika altyapısına dayanan ve kurumsal ortamlarda yaygın olarak kullanılan bir diğer güvenli e-posta standardıdır \cite{RamsdellTurner}.

RSA algoritması, Rivest, Shamir ve Adleman tarafından 1978 yılında tanıtılmış; asimetrik şifreleme ve dijital imzanın temelini oluşturmuştur \cite{RSA1978}. PKCS \#1 v2.2 standardı (RFC 8017), RSA-PSS imza şemasını ve RSA-OAEP şifreleme şemasını tanımlamaktadır \cite{RFC8017}. AES-256-GCM, NIST tarafından FIPS 197 ve SP 800-38D standartlarıyla belgelenmiş, kimlik doğrulamalı şifreleme sağlayan AEAD modudur \cite{FIPS197, SP80038D}. SHA-256 ise FIPS 180-4 standardıyla tanımlanan 256 bit çıktı üreten güvenli özet algoritmasıdır \cite{FIPS1804}.

Kriptografik sistemlerin görselleştirilmesi üzerine yapılan araştırmalar, animasyonlu gösterimlerin kriptografik kavramların anlaşılmasını önemli ölçüde kolaylaştırdığını ortaya koymaktadır \cite{CryptoVis}. Bu çalışmada geliştirilen animasyon modülleri, literatürdeki bu yaklaşımdan ilham almaktadır.

\textbf{Çalışmanın Amacı}

Bu çalışmasının amaçları şu şekilde sıralanabilir:

\begin{itemize}
  	\item Güvenli e-posta iletişiminin dört temel güvenlik özelliğini (gizlilik, bütünlük, kimlik doğrulama ve inkar edememe) karşılayan, endüstriyel kriptografik standartlara uygun bir hibrit şifreleme sistemi tasarlamak ve gerçekleştirmek.

  	\item AES-256-GCM, RSA-2048 PSS/OAEP ve SHA-256 algoritmalarının iç işleyişini; 14 turlu AES durum matrisi animasyonu, Genişletilmiş Öklid Algoritması adım animasyonu ve SHA-256 sıkıştırma devresi görselleştirmesi aracılığıyla pedagojik biçimde sunmak.

  	\item Geliştirilen sistemi 26 birim testi ile doğrulayarak yazılım kalitesini ve kriptografik doğruluğu güvence altına almak.
\end{itemize}

\textbf{Tezin Organizasyonu}

Bu çalışmanın düzenlemesi şu şekildedir.

1. Bölümde, e-posta güvenliğinin önemi, çalışmanın amacı ve literatürdeki ilgili çalışmalar ele alınmaktadır.

2. Bölümde, sistemde kullanılan kriptografik primitifler (SHA-256, RSA-2048, AES-256-GCM) teorik altyapısıyla birlikte açıklanmaktadır.

3. Bölümde, sistemin mimari tasarımı, modüler yapısı, GUI gerçekleştirimi ve animasyon modülleri ayrıntılı biçimde aktarılmaktadır.

4. Bölümde, birim test sonuçları ve güvenlik analizi sunulmakta; çalışmanın bulguları, sınırlılıkları ve öneriler tartışılmaktadır.

